\scanpage{650}

\[
(60)\qquad
\left\{\begin{aligned}
\Gamma_{\rho,\sigma}&\simeq
\Gamma_{\rho,\sigma}^{\unsure{\natural|\sharp}}\cdot
\underbrace{S\Gamma_{\rho,\sigma}}_{SZ_{\rho}\times SZ_{\sigma}=L_{\rho}\times L_{\sigma}}
\ ;\quad
\Gamma_{\rho,\sigma}^{\unsure{\natural|\sharp}}\cap S\Gamma_{\rho,\sigma}=\mu_{\rho}
\\
\Gamma_{\rho,\sigma}/\mu_{\rho}&\simeq
\left(\underbrace{\Gamma_{\rho,\sigma}^{\unsure{\natural|\sharp}}/\mu_{\rho}}
_{\gamma_{\rho,\sigma}\simeq\widetilde{\Gamma}_{\mathbb{Q}}}\right)
\cdot\left(L_{\rho}/\mu_{\rho}\times L_{\sigma}\right).
\end{aligned}\right.
\]

To say what becomes of the remarkable subgroups
\[
M(0),\quad
\begin{array}{c}
M[0],\\
\inserted{\Vert}\\[-4pt]
\inserted{M^{\mathcal{G}}[0]}\\
\inserted{\Vert}\\[-4pt]
\inserted{M(\unsure{\omega|0}).L_{0}^{\mathcal{G}}}
\end{array}
\quad
\begin{array}{c}
M_{\unsure{a}}^{!}[0],\\
\inserted{\Vert}\\[-4pt]
\inserted{M^{\unsure{\Pi}}[0]}\\
\inserted{\Vert}\\[-4pt]
\inserted{M(\unsure{\omega|0}).L_{0}^{\Pi}}
\end{array}
\quad
\begin{array}{c}
M_{0}^{\unsure{!}}[0]\\
\inserted{\Vert}\\[-4pt]
\inserted{M^{\Pi_{0}}\unsure{[0]}}\\
\inserted{\Vert}\\[-4pt]
\deleted{\unsure{M(\omega)L_{0}|M(0)L_{0}}}\\
\inserted{M(\unsure{\omega|0})L_{0}^{\Pi_{0}}}
\end{array}
\]
\[
\begin{array}{c}
\unsure{\{M(\tfrac12)\}|M(\tfrac12)}\\
\inserted{\Vert}\\[-4pt]
\inserted{Z_{\sigma0}^{!}}
\end{array},\quad
M'(-\bar{\jmath}),\quad M'''(-\bar{\jmath}),\quad N_{\rho}^{?}
\]
of $M$ under the homomorphism
\diagram{
\begin{tikzcd}[ampersand replacement=\&, column sep=large]
M \arrow[r,"\sim"] \&
\mathcal{M}_{\rho,\sigma}^{\unsure{\natural|\sharp}}/\mu_{\rho} \arrow[r,hook] \&
\mathcal{M}_{\rho,\sigma},
\end{tikzcd}}
we would first have to introduce the appropriate groups in the general $\mathcal{M}_{\rho,\sigma}$.
We are therefore going to return to these groups now.

\mspage{649}

\scanpage{651}
\mspage{628}

\uline{VI}\quad The general groups $\mathcal{M}_{\rho,\sigma}$:
recapitulation, refinement.

$1^\circ)$ Data. These are those of \uline{V} $1^\circ)$ (page~577).
In addition to $\mathcal{G}$ \inserted{(profinite)} and to
\diagram{\begin{tikzcd}[ampersand replacement=\&]
\mathrm{Gl}(2,\mathbb{Z})^{\wedge} \arrow[r] \& \mathcal{G}
\end{tikzcd}}
a continuous surjective homomorphism, i.e. generators $\rho,\sigma,\tau$
(or $\rho,\sigma,\tau'$) of $\mathcal{G}$, we also give ourselves an open subgroup
\[
(1)\qquad \mathcal{G}_{\unsure{0|a}}^{\unsure{\natural|\sharp}}
\subset\mathcal{G}^{+},\quad
\text{invariant in }\mathcal{G}^{\unsure{*|+}}.
\]

In the universal case $\mathcal{G}=\mathrm{Gl}(2,\mathbb{Z})^{\wedge}$,
the datum of $\mathcal{G}_{\unsure{0|a}}^{\unsure{\natural|\sharp}}$
may correspond to the point of view where we look at the subgroups
$M^{\unsure{\natural|\sharp}}$ of $M$
\deleted{which}\inserted{\unsure{formed|formed} of the $u$ which}
\unsure{normalize}\deleted{\unsure{ateur of|ent the}}
$\mathcal{G}^{\unsure{\natural|\sharp}}$
\inserted{and which operate \unsure{trivially} in
$\mathcal{G}^{+}/\mathcal{G}^{\unsure{\natural|\sharp}}$ --- this is}
\deleted{--- which is \unsure{necessarily} open, and defines}
\inserted{an open subgroup, which defines}
in $\widetilde{\Gamma}_{\mathbb{Q}}=M/\mathcal{G}^{+}$
an open image subgroup
$\widetilde{\Gamma}_{\mathbb{Q}}^{\unsure{\natural|\sharp}}$,
whose representations in suitable profinite groups we may propose to study.
An interesting case will be the one where
\[
\mathcal{G}^{\unsure{\natural|\sharp}}=\Pi_{0},\quad
\text{thus }M^{\unsure{\natural|\sharp}}=M^{\prime\prime!},
\]
or
\[
\mathcal{G}^{\unsure{\natural|\sharp}}=\Pi,\quad
\text{thus }M^{\unsure{\natural|\sharp}}=M^{!}.
\]
If
\[
\mathcal{G}^{\unsure{\natural|\sharp}}=\mathcal{G},\quad
\text{thus }M^{\unsure{\natural|\sharp}}=M,
\]
we recover the preceding constructions.
\deleted{\unsure{We can}}

\mspage{650}

\scanpage{652}

An interesting case also is the one where $\unsure{\mathcal{G}}$ is the group
of adelic \inserted{integral} points of a group scheme over $\mathbb{Z}$
(in this case, it is true that $\rho$ and $\sigma,\tau$ generate
\unsure{perhaps|not then} $\mathcal{G}^{\unsure{+|*}}$ only modulo a finite group
i.e. they generate \unsure{perhaps} only an open subgroup.
In this case, we should rather \uline{\unsure{paraphrase}} the constructions
that follow in the scheme-theoretic case
I shall come back to this later on. An interesting choice of a
$\mathcal{G}^{\unsure{\natural|\sharp}}$ is then the
\uline{identity component} of $\mathcal{G}$.

I hesitate to introduce notations
$Z_{\rho,\sigma}^{\unsure{\natural|\sharp}},\mathcal{M}_{\rho,\sigma}^{\unsure{\natural|\sharp}}$
etc --- so as not to overload notations that are already heavy. So I prefer to
use the simpler notations $Z_{\rho,\sigma}$, etc --- it being understood
that in these constructions, the group $\mathcal{G}^{\unsure{\natural|\sharp}}$
enters --- with the option of reintroducing the notation $\unsure{\natural|\sharp}$
if need be.
We put
\[
(3)\qquad
\mathcal{G}_{\tau}^{\unsure{\natural|\sharp}}
=\overbrace{\{1,\tau\}}^{\mu_{\tau}}\,.
\mathcal{G}^{\unsure{\natural|\sharp}}\subset\mathcal{G}.
\]

\mspage{651}

\scanpage{653}
\mspage{629}

$2^\circ)$ \quad $Z_{\rho,\sigma},\ \gamma_{\rho,\sigma}$.

Let
\[
(4)\qquad
\begin{array}{l}
 Z_{\rho,\sigma}\subset
 \operatorname{Aut}(\mathcal{G}^{+})\times
 \widehat{\mathbb{Z}}^{*}\times\mu\times\mu\\
 \phantom{Z_{\rho,\sigma}}\Vert\,\mathrm{dfn}\\
 \left\{(u,\mu,\varepsilon_2,\varepsilon_3)\ \left|
 \begin{array}{ll}
 \multicolumn{2}{l}{\deleted{\unsure{u\ \text{normalizes}\ \mathcal{G}^{\natural}\ 
 \text{and induces the identity in}\ \mathcal{G}^{+}/\mathcal{G}^{\natural}}}}\\
 a)&\exists\alpha\in\mathcal{G}_{\tau}^{\unsure{\natural|\sharp}},
 \ \text{with }\varepsilon_{\mathcal{G}}(\alpha)=\varepsilon_3,
 \quad u(\rho)=\alpha(\rho)\\
 b)&\exists\beta\in\mathcal{G}_{\tau}^{\unsure{\natural|\sharp}},
 \ \text{with }\varepsilon_{\mathcal{G}}(\beta)=\varepsilon_2,
 \quad u(\sigma)=\beta(\sigma)\\
 c)&u(\varepsilon_0)=\varepsilon_0^\mu\\
 d)&u\ \text{normalizes }\mathcal{G}^{\unsure{\natural|\sharp}}
 \inserted{\ \text{\uline{and} induces the identity on }
 (\mathcal{G}^{+}/\mathcal{G}^{\unsure{\natural|\sharp}})/\operatorname{Im}\Pi}
 \end{array}\right.\right\}.
\end{array}
\]
\margin{\unsure{cf. page} \unsure{635} \unsure{continuation}
\deleted{\unsure{a weaker condition}}
\unsure{this condition}.}
\margin{\deleted{The last condition \unsure{added} in d)
(\unsure{in principle}), is a \unsure{consequence} of a)b) if
$\varepsilon_2\varepsilon_3=1$
(\unsure{in which case} $u|(\mathcal{G}^{+}/\mathcal{G}^{\unsure{\natural|\sharp}})$
\deleted{the identity} \unsure{induced by} $\varepsilon[\tau]$,
\unsure{where} $\varepsilon=\varepsilon_2=\varepsilon_3\in\{\pm1\}$)}.}

\uline{Proposition.} \quad $Z_{\rho,\sigma}$ is a \uline{subgroup}
of the product group.

\uline{Proof.} The conditions c), d) define \unsure{obviously}
\deleted{a} \inserted{each} subgroup.
\unsure{Given that}
\[
\tau(\sigma)=-\unsure{\sigma|\varepsilon},
\]
condition b) also means that
\[
\unsure{(5^{*})}\qquad
b')\quad u(\sigma)\ \text{is }\mathcal{G}^{\unsure{\natural|\sharp}}
\text{-conjugate to : }\
\deleted{\unsure{\varepsilon_2}}\sigma^{\varepsilon_2}
\bigl(\deleted{\unsure{\varepsilon_2\sigma}}\inserted{=\varepsilon_2\sigma}\bigr)
\]
and in this form it is clear that condition c')
defines a subgroup. I do not think that condition a)
by itself defines a subgroup. But if
$\unsure{\mathfrak{g}_{(b)}|G_{(b)}}$
\deleted{$\unsure{\mathfrak{g}_{(a,b)}}$}
denotes the sub\unsure{group} of the product defined by these conditions
\deleted{\unsure{a)}}\inserted{b)}, we have \unsure{to show} that
$\unsure{\mathfrak{g}|G}_{\deleted{\unsure{(a,b,c)}}\inserted{(a,b)}}$
is a subgroup of the \deleted{product}
$\unsure{\mathfrak{g}|G}_{\deleted{\unsure{(a,b)}}\inserted{(b)}}$
--- which will prove the proposition.
Let $\uline{u},\uline{u}'\in\unsure{\mathfrak{g}_{(b)}|G_{(b)}}$,
corresponding to
\inserted{$(u,\mu,\varepsilon_2,\varepsilon_3)$ and
$(u',\mu',\varepsilon'_2,\varepsilon'_3)$, and to}
$\alpha,\alpha'\in\mathcal{G}^{\unsure{\natural|\sharp}}_\tau$,
\unsure{so} that we have

\mspage{652}

\scanpage{654}

\[
(6)\qquad
u(\rho)=\alpha(\rho),\quad u'(\rho)=\alpha'(\rho),
\inserted{\quad\text{with }\begin{array}{l}
\varepsilon_{\mathcal{G}}(\alpha)=\varepsilon_3,\\
\alpha_{\mathcal{G}}(\alpha')=\varepsilon'_3.
\end{array}}
\]
\unsure{Let us prove} \deleted{$\unsure{\varepsilon_3}$} that
$u'u\in\unsure{\mathfrak{g}_{(a)(b)}|G_{(a,b)}}$
--- i.e. that it also satisfies \unsure{(b)|(a)}.
If $\varepsilon_3=1$ i.e.
$\alpha\in\mathcal{G}^{\unsure{\natural|\sharp}}
\inserted{\subset\mathcal{G}^{+}}$,
$u'(\alpha)$ is defined and we have
\[
\begin{aligned}
(u'u)(\rho)&=u'(u(\rho))=u'(\alpha(\rho))
=u'(\alpha)(u'(\rho))=u'(\alpha)(\alpha'(\rho))\\
&=(u'(\alpha)\alpha')(\rho)=\alpha''(\rho),
\end{aligned}
\]
with
\[
(7)\qquad
\begin{array}{l|l}
\alpha''=u'(\alpha)\alpha'\in
\mathcal{G}^{\unsure{\natural|\sharp}}_\tau&
\text{Case }\varepsilon_3=1\\
\varepsilon_{\mathcal{G}}(\alpha'')
=\varepsilon_{\mathcal{G}}(\alpha')
=\varepsilon'_3=\varepsilon_3\varepsilon'_3&
\end{array}
\]
(since $u'(\alpha)\in\mathcal{G}^{\unsure{\natural|\sharp}}$
i.e. $\varepsilon_{\mathcal{G}}(u'(\alpha))=1$, and $\varepsilon_3=1$.

In this case we \unsure{win}. If $\varepsilon_3=-1$, we write
\[
\alpha=\alpha_0\tau\qquad
\alpha_0\in\mathcal{G}^{\unsure{\natural|\sharp}},
\]
and we have
\[
u(\rho)=(\alpha_0\tau)(\rho)=\alpha_0(\tau(\rho))
=\alpha_0\bigl(\sigma^{\unsure{-1}}(\rho^{-1})\bigr)
=(\alpha_0\sigma)(\rho^{-1}),
\]
where now $\alpha_0\sigma\in\mathcal{G}^{+}$, hence
$u'(\alpha_0\sigma)$ is defined, and we will have
\[
\unsure{(u'u)(\rho)}
=u'(u(\rho))
=u'(\alpha_0\sigma)
\bigl(\underbrace{u'(\rho^{-1})}_{\alpha'(\rho^{-1})}\bigr)
=\bigl(u'(\alpha_0\sigma)\alpha'\bigr)(\rho^{-1})
\]
and since
\[
\deleted{\unsure{\rho^{-1}}}\qquad
\rho^{-1}=\sigma\tau(\rho)
\]
we find
\[
(8)\qquad
(u'u)(\rho)=
\deleted{\unsure{u'(\alpha_0\sigma)\alpha'\sigma\tau}}
\alpha''(\rho)
\]
with
\[
\left\{
\begin{array}{l}
\alpha''=\varepsilon
\overbrace{u'(\alpha_0)}^{\inserted{\in\mathcal{G}^{\unsure{\natural|\sharp}}}}
\overbrace{u'(\sigma)}^{\inserted{\in\mathcal{G}^{+}}}
\overbrace{\alpha'}^{\inserted{\in\mathcal{G}^{\unsure{\natural|\sharp}}_\tau}}
\underset{\deleted{\unsure{\tau'}}}{\sigma\tau}
\qquad\text{where}\\
\varepsilon\in\mu=\{\pm1\}
\unsure{\hookrightarrow|\Rightarrow}\operatorname{Centre}\mathcal{G}^{+}
\end{array}
\right.
\]
\deleted{\unsure{is}} is chosen so that $\alpha''$ (if possible)
is in $\mathcal{G}^{\unsure{\natural|\sharp}}_\tau$.
Let us note that all the factors of this product are in
$\unsure{\mathcal{G}^{\natural}_\tau\mathcal{G}^{+}|
\mathcal{G}^{\natural}_\tau/\mathcal{G}^{+}}$,
\unsure{\uline{in}} any case $\alpha''\in\mathcal{G}$.

\mspage{653}